\documentclass[11pt]{article}
\usepackage[utf8]{inputenc}
\usepackage[T1]{fontenc}
\usepackage[italian]{babel}
\usepackage{geometry}
\geometry{a4paper, margin=1in}
\usepackage{xcolor}
\usepackage{listings}
\usepackage{hyperref}

\usepackage{tcolorbox}
\tcbuselibrary{breakable}
\begin{document}

\begin{tcolorbox}[colback=blue!5!white,colframe=blue!75!black,title=\textbf{DOMANDA}, breakable]
3) Mettere il server in ascolto sulla porta 100 e osservare cosa succede
\begin{itemize}
    \item Bisogna modificare anche il client? Dove?
    \item Per chi usa il proprio PC con Linux o una virtual machine Linux, lanciare il server con il comando ``\texttt{sudo ./serverUDP}'' e osservare cosa cambia
\end{itemize}
\end{tcolorbox}

\begin{tcolorbox}[colback=green!5!white,colframe=green!75!black,title=\textbf{RISPOSTA}, breakable]
Ecco le risposte alle tue domande:

\begin{itemize}
    \item \textbf{Bisogna modificare anche il client? Dove?}
    
    Sì, è indispensabile modificare anche il client. Affinché il client riesca a contattare il server, deve inviare il messaggio alla porta corretta. Nel file \texttt{clientUDP.c} è necessario modificare il numero di porta di destinazione passato come ultimo argomento alla funzione \texttt{UDPSend} (portandolo da 10000 a 100). Nel codice aggiornato, questo si traduce nell'impostare la macro \texttt{SERVERPORT} a 100 anche per il client.
    
    \item \textbf{Per chi usa il proprio PC con Linux... lanciare il server con ``\texttt{sudo ./serverUDP}'' e osservare cosa cambia}
    
    Sui sistemi Linux (e Unix in generale), le porte di rete numerate da 0 a 1023 sono definite "privileged ports" (o well-known ports). Se tenti di eseguire \texttt{./serverUDP} in ascolto sulla porta 100 come utente normale, il sistema operativo impedirà l'operazione. La funzione di sistema per riservare la porta (\texttt{bind()}) fallirà restituendo un errore di "Permission denied". Di conseguenza, il server non si avvierà.
    
    Utilizzando invece \texttt{sudo ./serverUDP}, il programma viene eseguito con i privilegi massimi (utente root), che possiede l'autorizzazione per vincolare i processi alle porte privilegiate (< 1024). In questo caso, il server si avvierà con successo e si metterà normalmente in ascolto sulla porta 100.
\end{itemize}
\end{tcolorbox}

\end{document}
